\documentclass[fontsize=15pt]{jlreq}

%プリアンブル
\usepackage{amsmath}
\usepackage{mathtools}
\pagestyle{empty}
\usepackage[top=15mm, bottom=20mm, left=20mm, right=20mm]{geometry}

\newcommand{\m}[1]{\raisebox{0.1em}{\textcircled{\scriptsize #1}}}



%本文
\begin{document}
\begin{center}
{\LARGE {振り返り}} \\
-物理\m{4}- %単元ごとに変えるか消す
\end{center}
\vspace{1em}
\begin{flushright}
名前：\underline{\hspace{5cm}}
\end{flushright}

% 問題部分
\begin{enumerate}
    \item 慣性の法則について次の語句を全て含めて説明せよ．\\
    \small 語句：力，静止，運動 \\
    回答：
    \vspace{5em} 

    \item 作用・反作用の法則の説明せよ．\\
    (\hspace{15cm})
    \vspace{2em} 

    \item 水平面の上に置かれた物体に作用する力について考える．物体にはたらく重力と，水平面が物体を押し返す垂直抗力の関係について，次の4つのうち正しいものを選べ．
    \begin{itemize}
      \noindent  A：2力の大きさは，物体の質量と水平面の質量が等しい場合にのみ等しくなる．\\
        B：2つの力は釣り合っている．\\
        C：2力はどちらも水平面に作用している．\\
        D：2力は作用・反作用の関係にある．
    \end{itemize}
    (\hspace{2cm})
    \vspace{2em} 

    \item 氷の上で静止している2人の人AとBがいる．AがBを前に押したとき，A，Bはそれぞれどのような運動をするか．\\
    (\hspace{15cm})
    \vspace{2em}

    \item 水圧と浮力はそれぞれ何に比例しているか．
    \begin{itemize}
        \item 水圧：
        \item 浮力：
    \end{itemize}
    \vspace{2em}  

\end{enumerate}
\end{document}
